% Part: many-valued-logic
% Chapter: syntax-and-semantics
% Section: connectives

\documentclass[../../../include/open-logic-section]{subfiles}

\begin{document}

\olfileid{mvl}{syn}{con}

\olsection{ભાષાઓ અને સંયોજકો}

શાસ્ત્રીય વિધાનાત્મક તર્કશાસ્ત્ર અને અન્ય ઘણાં
તર્કશાસ્ત્રો \emph{વિધાનાત્મક અચલો} તથા
\emph{સંયોજકો}નો ગણ વાપરે છે. દાખલા તરીકે,
નીચેનાંને આપણે મૂળભૂત માનીએ છીએ:
\begin{enumerate}
\tagitem{prvFalse}{!!{falsity} માટેનો વિધાનાત્મક અચલ~$\lfalse$.}{}
\tagitem{prvTrue}{!!{truth} માટેનો વિધાનાત્મક અચલ~$\ltrue$.}{}
\item તાર્કિક સંયોજકો:
  \startycommalist
  \iftag{prvNot}{\ycomma $\lnot$ (નિષેધ)}{}%
  \iftag{prvAnd}{\ycomma $\land$ (સંયોજન)}{}%
  \iftag{prvOr}{\ycomma $\lor$ (વિયોજન)}{}%
  \iftag{prvIf}{\ycomma $\lif$ (!!{conditional})}{}%
  \iftag{prvIff}{\ycomma $\liff$ (!!{biconditional})}{}%
\end{enumerate}
\iftag{defNot,defOr,defAnd,defIf,defIff,defTrue,defFalse,defEx,defAll}{%
ઉપરના મૂળભૂત સંયોજકો ઉપરાંત આપણે સંક્ષેપો તરીકે
વ્યાખ્યાયિત સંકેતો પણ વાપરીએ છીએ, જેમ કે
\startycommalist
  \iftag{defNot}{\ycomma $\lnot$ (નિષેધ)}{}%
  \iftag{defAnd}{\ycomma $\land$ (સંયોજન)}{}%
  \iftag{defOr}{\ycomma $\lor$ (વિયોજન)}{}%
  \iftag{defIf}{\ycomma $\lif$ (!!{conditional})}{}%
  \iftag{defIff}{\ycomma $\liff$ (!!{biconditional})}{}%
  \iftag{defFalse}{\ycomma $\lfalse$ (!!{falsity})}{}%
  \iftag{defTrue}{\ycomma $\ltrue$ (!!{truth})}.}{}

બહુમૂલ્ય તર્કશાસ્ત્રોમાં પણ આ જ સંયોજકો વપરાય છે.
જોકે, એક જ તર્કશાસ્ત્રમાં, માનો કે સંયોજનનાં જુદાં
સ્વરૂપો સમાવવાં ઘણી વાર ઉપયોગી છે; તેમને અલગ
ઓળખાવવા જુદા સંકેતો જોઈએ. કેટલાક બહુમૂલ્ય
તર્કશાસ્ત્રોમાં શાસ્ત્રીય તર્કશાસ્ત્રમાં સમકક્ષ ન હોય
એવા સંયોજકો પણ છે. તેથી આપણે સામાન્ય કરતાં
થોડી વધુ વ્યાપક રીત લઈશું.

\begin{defn}
  \emph{વિધાનાત્મક ભાષા} એટલે સંયોજકોનો ગણ~$\Lang L$.
  દરેક સંયોજક~$\star$ની \emph{આર્ગ્યુમેન્ટ-સંખ્યા}
  હોય છે; આર્ગ્યુમેન્ટ-સંખ્યા~$n$ ધરાવતો સંયોજક
  \emph{$n$-સ્થાનીય} કહેવાય છે. આર્ગ્યુમેન્ટ-સંખ્યા~$0$
  ધરાવતા સંયોજકો \emph{અચલો} પણ કહેવાય છે;
  આર્ગ્યુમેન્ટ-સંખ્યા~$1$ ધરાવતાં સંયોજકોને
  \emph{એકસ્થાનીય} અને આર્ગ્યુમેન્ટ-સંખ્યા~$2$
  ધરાવતાંને \emph{દ્વિસ્થાનીય} કહે છે.
\end{defn}

\begin{ex}
  વિધાનાત્મક તર્કશાસ્ત્રની પ્રમાણભૂત ભાષા~$\Lang L_0$માં
  નીચેના સંયોજકો છે (કૌંસમાં તેમની આર્ગ્યુમેન્ટ-સંખ્યા છે):
  $\lfalse$~($0$),
  $\lnot$~($1$),
  $\land$~($2$),
  $\lor$~($2$),
  $\lif$~($2$). આપણે વિચારવાના મોટાભાગના
  તર્કશાસ્ત્રો આ ભાષા વાપરશે. કેટલાક તર્કશાસ્ત્રો
  રૂઢિ અનુસાર અમુક સંયોજકો માટે જુદા સંકેતો
  વાપરે છે. દાખલા તરીકે, ગુણન તર્કશાસ્ત્રમાં
  સંયોજનનો સંકેત ઘણી વાર $\land$ને બદલે~$\odot$
  હોય છે. ક્યારેક નવો કારક ઉમેરવો અનુકૂળ હોય છે;
  દાખલા તરીકે, નિર્ધારિતતા કારક~$\triangle$
  ($1$-સ્થાનીય).
\end{ex}

\end{document}
